``我们已经把姓名删掉了。''

这句话在很多数据发布流程里被当作隐私工作的终点，而它其实连起点都算不上。删掉姓名之后剩下的字段仍然能组合出唯一标识，补救的办法是继续删，继续删的终点是一份没有信息的数据——两端各有一个不能同时满足的要求，中间没有落脚点。

这个单元讲的是绕开这条死路的一次问法改变。不问``输出里还剩多少个人信息''，改问``某一个人在与不在数据集里，输出的分布差多少''。后者可以精确定义，可以量化，而且不需要假设攻击者手上有什么外部数据。

改变问法要付出的东西在定义写出来的那一刻就注定了：这个条件是一个关于随机机制的全称命题，要让它成立必须注入噪声，而噪声的量由一条记录能把答案推多远决定，与数据规模无关。更要紧的是，这份保证会随着查询次数被消耗，直到一分不剩。三篇分别处理这三件事。

\hyperref[sec:17-Constrained-Guarantees/01-Differential-Privacy/01]{``匿名化为什么不是隐私''}从一次真实的重识别攻击出发，说明为什么删列这条路没有终点，然后给出 \(\varepsilon\)-差分隐私的定义，重点拆开它那两个全称量词——定义的全部强度都在那里，同时它是关于机制的性质而不是关于某次输出的性质。

\hyperref[sec:17-Constrained-Guarantees/01-Differential-Privacy/02]{``加多少噪声由敏感度决定''}回答工程上最直接的问题。答案是敏感度，不是数据规模、不是答案量级。Laplace 机制的证明只有三行，而它的推论很硬：计数查询的敏感度恒为 1，于是同样的噪声在百万级数据上是尘埃，在百人级数据上会毁掉结果。

\hyperref[sec:17-Constrained-Guarantees/01-Differential-Privacy/03]{``预算会被花光''}把单次发布的账扩展到长期运行的系统。组合定理说明 \(\varepsilon\) 是相加的，后处理不变性说明加噪之后的数据可以随意使用，两条合起来决定了一个隐私系统该怎么设计。DP-SGD 在这里出现，它把深度学习的训练循环整个装进这套记账规则。

读完之后值得带走的判断有两条。一条是隐私不是标签而是预算，任何声称``本系统满足差分隐私''却不说累计 \(\varepsilon\) 的表述都是不完整的。另一条更一般：当两个目标之间的换算关系可以被写成不等式，讨论就从``做得够不够好''转成了``工作点选在哪里''，而选点是一个需要有人负责的判断，不是一个可以外包给算法的技术细节。
